\documentclass[12pt,oneside,a4paper]{article}

% ---------------------------------------------------------
% PACOTES BÁSICOS
% ---------------------------------------------------------
\usepackage[T1]{fontenc}
\usepackage[brazil]{babel}
\usepackage{times}
\usepackage{setspace}
\usepackage{geometry}
\usepackage{graphicx}
\usepackage[export]{adjustbox}
\usepackage{float}
\usepackage{listings}
\usepackage{xcolor}
\usepackage{caption}
\usepackage{booktabs}
\usepackage{hyperref}
\usepackage{xurl}
\usepackage{microtype}
\usepackage{ragged2e}
\usepackage{seqsplit}
\usepackage{enumitem}
\usepackage{amsmath, amssymb}
\usepackage{indentfirst}
\usepackage{array}
\usepackage{tabularx}
\usepackage[section]{placeins}
\usepackage[num]{abntex2cite}

\graphicspath{{../arquitetura/diagramas/imagens/}{imagens/}}

\lstset{
    language=C++,
    basicstyle=\ttfamily\footnotesize,
    keywordstyle=\color{blue}\bfseries,
    backgroundcolor=\color{black!5},
    stringstyle=\color{red!70!black},
    commentstyle=\color{gray!90!black},
    showstringspaces=false,
    breaklines=true,
    breakatwhitespace=false,
    columns=fullflexible,
    keepspaces=true,
    frame=single,
    numbers=left,
    numberstyle=\tiny\color{gray},
    tabsize=2,
    extendedchars=true,
    literate={á}{{\'a}}1 {ã}{{\~a}}1 {é}{{\'e}}1 {ç}{{\c{c}}}1 {ó}{{\'o}}1 {í}{{\'i}}1
}

\renewcommand{\lstlistingname}{Código}
\renewcommand{\lstlistlistingname}{Lista de Códigos}

% ---------------------------------------------------------
% FORMATAÇÃO
% ---------------------------------------------------------
\geometry{a4paper, left=3cm, right=2cm, top=2.5cm, bottom=2.5cm}
\setstretch{1.5}
\captionsetup{labelfont=bf, labelsep=endash, justification=centering}

\justifying
\sloppy
\emergencystretch=3em
\setlength{\parindent}{1.25cm}
\setlength{\parskip}{0pt}

% Permite quebra segura de termos técnicos longos em fonte monoespaçada
\let\OldTexttt\texttt
\renewcommand{\texttt}[1]{\begingroup\ttfamily\seqsplit{#1}\endgroup}

% Listas dentro das margens
\setlist[itemize]{
  leftmargin=1.15cm,
  itemsep=0.25em,
  topsep=0.4em,
  parsep=0pt
}

\setlist[enumerate]{
  leftmargin=1.15cm,
  itemsep=0.35em,
  topsep=0.4em,
  parsep=0pt
}

% Espaçamento controlado ao redor de figuras e tabelas
\setlength{\textfloatsep}{10pt plus 2pt minus 2pt}
\setlength{\floatsep}{10pt plus 2pt minus 2pt}
\setlength{\intextsep}{10pt plus 2pt minus 2pt}
\setlength{\abovecaptionskip}{6pt}
\setlength{\belowcaptionskip}{4pt}

% ---------------------------------------------------------
% COMANDOS PARA FIGURAS EM RETRATO, DENTRO DAS MARGENS
% ---------------------------------------------------------
\newcommand{\FiguraRetrato}[3]{
\begin{figure}[!htbp]
\centering
\includegraphics[
  max width=0.88\textwidth,
  max height=0.62\textheight,
  keepaspectratio
]{#1}
\caption{#2}
\label{#3}
\end{figure}
\FloatBarrier
}

\newcommand{\FiguraRetratoGrande}[3]{
\begin{figure}[!htbp]
\centering
\includegraphics[
  max width=0.90\textwidth,
  max height=0.68\textheight,
  keepaspectratio
]{#1}
\caption{#2}
\label{#3}
\end{figure}
\FloatBarrier
}

\newcommand{\FiguraManual}[3]{
\begin{figure}[!htbp]
\centering
\includegraphics[
  max width=0.88\textwidth,
  max height=0.78\textheight,
  keepaspectratio
]{#1}
\caption{#2}
\label{#3}
\end{figure}
\FloatBarrier
}

% ---------------------------------------------------------
% DADOS DO TRABALHO
% ---------------------------------------------------------
\newcommand{\Instituicao}{INSTITUTO FEDERAL DE EDUCAÇÃO, CIÊNCIA E TECNOLOGIA DO PIAUÍ}
\newcommand{\Campus}{Campus Piripiri}
\newcommand{\Curso}{Tecnólogo em Análise e Desenvolvimento de Sistemas}
\newcommand{\Titulo}{Sistema Embarcado para \textit{Feedback} Háptico em \textit{Sim Racing}}
\newcommand{\Autor}{Cairon Ferreira Prado}
\newcommand{\Orientador}{[Preencher nome do(a) orientador(a)]}
\newcommand{\CoOrientador}{}
\newcommand{\Cidade}{Piripiri -- PI}
\newcommand{\Ano}{2026}

\begin{document}

% ================= CAPA =====================
\begin{titlepage}
\begin{center}
% \includegraphics[width=10cm]{imagens/Logo-IFPI-Vertical.png} \\[0.3cm]
\textbf{\Instituicao} \\
\textbf{\Campus} \\
\textbf{\Curso} \\[3cm]

\textbf{\Autor} \\[2.5cm]
\textbf{\Titulo} \\[6cm]

\Cidade\\
\Ano
\end{center}
\end{titlepage}

% ================= FOLHA DE ROSTO =====================
\begin{titlepage}
\begin{center}
\textbf{\Autor} \\[3cm]

\textbf{\Titulo} \\[2cm]

\begin{flushright}
\begin{minipage}{0.6\textwidth}
Relatório Técnico de Software apresentado ao curso de \Curso{}, como requisito parcial para a conclusão do Trabalho de Conclusão de Curso (TCC). \\ \\
Orientador(a): \Orientador
\ifx\CoOrientador\empty\else \\ Coorientador(a): \CoOrientador \fi
\end{minipage}
\end{flushright}

\vfill
\Cidade \\
\Ano
\end{center}
\end{titlepage}

% ================= RESUMO =====================
\newpage
\section*{Resumo}
Este relatório técnico apresenta o desenvolvimento e a validação de um firmware
embarcado para \texttt{ESP32} voltado a converter telemetria de simuladores de
corrida em \textit{feedback} físico e visual. O sistema recebe pacotes \texttt{UDP},
decodifica o formato \textit{Data Out} atualmente implementado para
\textit{Forza Motorsport 7}, atualiza instrumentos físicos de velocidade,
rotação do motor, combustível e temperatura média dos pneus, e exibe marcha,
posição, volta, conectividade e estados operacionais em um display
\texttt{TFT}. A solução foi estruturada em camadas, com influência de
\textit{Ports and Adapters}, para reduzir acoplamento, melhorar a testabilidade e
permitir evolução incremental. O projeto também incorpora portal assíncrono de
configuração \texttt{Wi-Fi}, serviço de telemetria com \textit{timeout} e recuperação de
sinal, e organização do painel por catálogos de instrumentos e \textit{layouts}. A
validação combinou testes unitários, integração \textit{host-side}, teste de componente,
\textit{smoke test} no microcontrolador e preparação de uma suíte de integração
\texttt{UDP} real em \textit{hardware}. Em verificação local realizada em 4 de maio de
2026, a suíte \texttt{test\_native} executou \texttt{68/68} casos com
\texttt{100\%} de aprovação; a documentação do projeto ainda registra um total
de \texttt{18} suítes e \texttt{73} casos definidos, incluindo \textit{smoke
test} embarcado aprovado e integração \texttt{UDP} real implementada, embora
dependente de uma bancada elétrica mais estável para fechamento experimental
completo. O resultado final é um produto mínimo viável (MVP) tecnicamente consistente, modular e apto à
demonstração acadêmica, com expansão futura prevista para compatibilidade
formal com outras plataformas e efeitos adicionais de vibração.

Palavras-chave: sistemas embarcados, ESP32, telemetria, \textit{sim racing}, firmware,
validação automatizada.

\newpage
\section*{Abstract}
This technical report presents the development and validation of an
\texttt{ESP32}-based embedded firmware designed to transform racing simulator
telemetry into physical and visual feedback. The system receives \texttt{UDP}
packets, decodes the currently supported \textit{Data Out} format from
\textit{Forza Motorsport 7}, updates physical gauges for speed, engine RPM,
fuel and average tire temperature, and renders gear, lap, position,
connectivity and operational status on a \texttt{TFT} display. The solution was
organized in layered form, influenced by \textit{Ports and Adapters}, in order
to reduce coupling, improve testability and support incremental evolution. The
project also includes an asynchronous \texttt{Wi-Fi} configuration portal, a
telemetry service with timeout and signal recovery, and a composition model
based on instrument and display catalogs. Validation combined unit tests,
host-side integration tests, component tests, an embedded smoke test and a real
\texttt{UDP} integration suite prepared for the microcontroller. In a local
verification performed on May 4, 2026, the \texttt{test\_native} suite executed
\texttt{68/68} test cases with full approval; project documentation also
records \texttt{18} suites and \texttt{73} defined cases in total, including a
successful embedded smoke test and a real \texttt{UDP} integration suite that
is still limited by electrical bench conditions for full experimental closure.
The final outcome is a technically consistent and modular MVP suitable for
academic demonstration, with future expansion planned for broader platform
compatibility and additional vibration effects.

Keywords: embedded systems, ESP32, telemetry, sim racing, firmware, automated
validation.

% ================= SUMÁRIO =====================
\newpage
\tableofcontents
\newpage

% ================= DESENVOLVIMENTO =====================
\section{Introdução}
\label{sec:introducao}

O contexto de \textit{sim racing} favorece projetos que ampliam a imersão do
jogador por meio de \textit{feedback} visual, sonoro e físico. Embora a telemetria dos
simuladores modernos seja rica em informações sobre velocidade, rotação do
motor, marcha, combustível, temperatura e posição em pista, essas informações
normalmente permanecem confinadas a interfaces digitais do próprio jogo ou a
aplicações complementares executadas em outras telas. No escopo deste TCC, o
problema técnico abordado consiste em transformar parte dessa telemetria em
comportamento observável de um painel físico, com baixo acoplamento ao \textit{hardware}
e com uma estratégia de validação compatível com as restrições de um protótipo
embarcado.

O repositório do projeto estabelece a fonte primária de verdade para este
relatório. Assim, o escopo final considerado não é o conjunto máximo de ideias
levantadas no pré-projeto, mas sim o que foi efetivamente implementado,
testado, documentado e versionado. Nessa versão final, o firmware materializa
um produto mínimo viável (MVP) centrado em cinco capacidades principais: configuração de rede por
portal \texttt{Wi-Fi}, recepção de telemetria via \texttt{UDP}, decodificação
do protocolo \textit{Forza Motorsport 7 (Data Out)}, acionamento de
instrumentos físicos e apresentação de estados e indicadores em \textit{display}
\texttt{TFT}.

Do ponto de vista arquitetural, a solução foi estruturada de forma a separar o
ponto de composição do sistema, a lógica de aplicação, os modelos de domínio,
as portas de abstração e os adaptadores concretos de infraestrutura. Essa
decisão permitiu exercitar grande parte do comportamento em ambiente \textit{host-side},
reduzindo a dependência do microcontrolador para cada iteração de
desenvolvimento. A organização incremental do projeto e a definição de
objetivos mensuráveis seguem uma abordagem metodológica compatível com a
proposta de pesquisa aplicada discutida por Wazlawick
\cite{wazlawick2017metodologia}.

\subsection{Objetivos}

\subsubsection{Geral}
Desenvolver e validar um sistema embarcado modular, fundamentado no uso do
microcontrolador \texttt{ESP32}, capaz de receber, processar e mapear dados de
telemetria provenientes de simuladores de corrida para um painel físico com
instrumentos analógicos e interface visual embarcada.

\subsubsection{Específicos}
\begin{itemize}
  \item Implementar a configuração e a reconexão de rede local por meio de
  portal \texttt{Wi-Fi} embarcado e reaproveitamento de credenciais salvas.
  \item Receber continuamente pacotes de telemetria via \texttt{UDP} na rede
  local utilizada pelo simulador.
  \item Decodificar o formato \textit{Forza Motorsport 7 (Data Out)} e extrair
  velocidade, rotação, marcha, volta, posição, combustível e temperaturas dos
  pneus.
  \item Traduzir os sinais de telemetria para atuação física em instrumentos de
  velocidade, \textit{RPM}, combustível e temperatura.
  \item Exibir, em display \texttt{TFT}, estados de conectividade, marcha,
  status operacional e \textit{layouts} de acompanhamento da telemetria.
  \item Estruturar o firmware de forma modular, favorecendo manutenção,
  testabilidade e expansão futura para novos instrumentos ou novos \textit{decoders}.
  \item Validar a solução por meio de testes automatizados \textit{host-side} e de
  execução controlada no \texttt{ESP32}.
\end{itemize}

\section{Tecnologias Envolvidas}

O projeto não envolve \textit{backend} web, banco de dados ou infraestrutura de
\textit{deploy} em nuvem. Trata-se de um firmware embarcado cujo núcleo técnico
se apoia em bibliotecas e ferramentas voltadas a redes locais, controle de
periféricos e validação automatizada.

\begin{table}[H]
\centering
\caption{Tecnologias efetivamente utilizadas no projeto}
\label{tab:tecnologias_usadas}
\footnotesize
\setlength{\tabcolsep}{4pt}
\renewcommand{\arraystretch}{1.25}
\begin{tabularx}{\textwidth}{
  >{\raggedright\arraybackslash}p{3.2cm}
  >{\justifying\arraybackslash}X
}
\toprule
\textbf{Tecnologia} & \textbf{Papel no projeto} \\
\midrule
\texttt{ESP32 Dev Module} & Plataforma de execução do firmware e integração com \texttt{Wi-Fi}, \texttt{GPIO} e periféricos do protótipo. \\
\texttt{C/C++} & Linguagem empregada na implementação do firmware e das suítes de teste. \\
\texttt{Arduino Framework} & Base de acesso ao ciclo \texttt{setup/loop}, temporização, pinos digitais, \texttt{Wi-Fi} e demais recursos do microcontrolador. \\
\texttt{PlatformIO} & Orquestração de \textit{build}, \textit{upload}, ambientes de teste e reprodução da validação. \\
\texttt{UDP} & Transporte de telemetria em rede local entre o simulador e o dispositivo embarcado. \\
\texttt{WiFi}, \texttt{WiFiUDP}, \texttt{WebServer} e \texttt{Preferences} & Conectividade local, recepção de pacotes, portal de configuração e persistência de credenciais. \\
\texttt{TFT\_eSPI} & Renderização dos \textit{layouts} no display \texttt{ST7789} de \texttt{240x320} pixels. \\
\texttt{TMC2208} & Driver empregado no acionamento dos motores associados aos ponteiros do painel. \\
\texttt{Unity} & \textit{Framework} de testes automatizados no host e no \texttt{ESP32}. \\
\texttt{PlantUML} e \texttt{Markdown} & Documentação arquitetural, modelagem técnica e rastreabilidade textual do TCC. \\
\bottomrule
\end{tabularx}
\end{table}

\section{Modelagem do Projeto}

\subsection{Levantamento de Requisitos}

O catálogo de requisitos do projeto foi mantido no repositório durante o
desenvolvimento e serviu como referência para o recorte do MVP final. Como o
escopo inicial era mais amplo do que a implementação entregue, a leitura dos
requisitos no presente relatório deve considerar três estados: implementado,
parcialmente contemplado e previsto como evolução futura.

\textbf{Requisitos Funcionais}
\begin{itemize}
  \item \texttt{RF01} -- conexão automática à rede local: implementado por meio
  do \texttt{WiFiConfigPortal}, com reaproveitamento de credenciais salvas e
  abertura de ponto de acesso quando a conexão falha.
  \item \texttt{RF02} -- comunicação com o jogo via \texttt{UDP}: implementado
  pelo adaptador \texttt{UdpReceiver} e pelo ambiente de telemetria do
  firmware.
  \item \texttt{RF03} -- interpretação dos dados recebidos: implementado pelo
  \texttt{Forza7Decoder} e pelo \texttt{TelemetryService}.
  \item \texttt{RF04} -- controle do indicador de velocidade: implementado por
  \texttt{SpeedGauge} e \texttt{GaugeMotorTmc2208}.
  \item \texttt{RF05} -- controle do indicador de rotação do motor:
  implementado por \texttt{RpmGauge} e \texttt{GaugeMotorTmc2208}.
  \item \texttt{RF06} -- controle do indicador de marcha: implementado na
  interface \texttt{TFT} por meio do \texttt{DisplayService} e dos layouts
  embarcados.
  \item \texttt{RF07} -- detecção automática de módulos conectados: apenas
  parcialmente contemplado. O software já é modular e aceita registro por
  catálogo, mas não há detecção física de \textit{hot-plug} no protótipo atual.
  \item \texttt{RF08} -- inicialização automática do sistema: implementado na
  classe \texttt{App}, que inicializa portal, telemetria, UI e instrumentos ao
  receber energia.
  \item \texttt{RF09} -- indicação de status operacional: implementado por meio
  do estado \texttt{UiStatus} e das telas de conexão, portal e telemetria.
  \item \texttt{RF10} a \texttt{RF13} -- compatibilidade formal com
  \texttt{Xbox Series}, \texttt{PlayStation 4}, \texttt{PlayStation 5} e
  \texttt{PC}: não há validação multiplataforma consolidada no repositório
  final; o \textit{decoder} efetivamente implementado é específico para
  \textit{Forza Motorsport 7}.
  \item \texttt{RF14} e \texttt{RF15} -- efeitos adicionais de vibração:
  permanecem fora do MVP entregue no firmware atual.
\end{itemize}

\textbf{Requisitos Não Funcionais}
\begin{itemize}
  \item \texttt{RNF01} -- facilidade de uso: parcialmente atendido pelo fluxo
  de portal \texttt{Wi-Fi}, que reduz configurações manuais e dispensa
  recompilação para troca de rede.
  \item \texttt{RNF02} -- baixa latência: orientou a escolha de \texttt{UDP} e
  o desenho do firmware, mas não possui medição formal fim a fim concluída nesta
  versão do trabalho.
  \item \texttt{RNF04} -- modularidade e expansibilidade: atendido pela
  organização em camadas, pelas portas e pelo registro por catálogos de \textit{layouts}
  e instrumentos.
  \item \texttt{RNF05}, \texttt{RNF11}, \texttt{RNF12} e \texttt{RNF13} --
  compatibilidade por plataforma: permanecem como objetivo de validação futura.
  \item \texttt{RNF06} -- robustez: fortemente contemplado nos testes por
  cenários de pacote inválido, descarte, \textit{timeout}, \textit{burst} e recuperação.
  \item \texttt{RNF08} -- escalabilidade funcional: parcialmente contemplado,
  pois a arquitetura permite novos \textit{decoders}, ainda que apenas um protocolo
  concreto esteja implementado.
  \item \texttt{RNF09} -- manutenibilidade: atendido pela organização do código,
  pela injeção de dependências e pela automação de testes.
  \item \texttt{RNF10} -- confiabilidade física: parcialmente contemplado. O
  firmware implementa limites lógicos e calibração dos instrumentos, mas a
  validação eletromecânica completa depende de fechamento experimental da
  bancada.
\end{itemize}

\subsection{Casos de Uso}

O catálogo analítico de casos de uso do TCC foi documentado no arquivo
\texttt{docs/analise/casos-de-uso.md}. Entretanto, considerando o repositório
final como referência principal, os casos de uso centrais efetivamente
materializados pelo firmware podem ser sintetizados conforme a
Tabela~\ref{tab:casos_uso}.

\begin{table}[H]
\centering
\caption{Casos de uso centrais do MVP final}
\label{tab:casos_uso}
\footnotesize
\setlength{\tabcolsep}{4pt}
\renewcommand{\arraystretch}{1.25}
\begin{tabularx}{\textwidth}{
  >{\raggedright\arraybackslash}p{1.7cm}
  >{\raggedright\arraybackslash}p{3.5cm}
  >{\justifying\arraybackslash}X
}
\toprule
\textbf{ID} & \textbf{Caso de uso} & \textbf{Materialização no repositório} \\
\midrule
\texttt{UC11} & Conectar automaticamente à rede local & O firmware tenta reutilizar credenciais salvas, inicia tentativa assíncrona de conexão e abre o portal \texttt{SimHub} quando necessário. \\
\texttt{UC12} & Receber telemetria via \texttt{UDP} & O adaptador \texttt{UdpReceiver} abre a porta \texttt{5300} e fornece pacotes ao serviço de telemetria. \\
\texttt{UC13} & Interpretar dados de telemetria & O \texttt{Forza7Decoder} converte o pacote bruto em um \texttt{TelemetryFrame} consumido pela aplicação. \\
\texttt{UC01} & Exibir velocidade no painel físico & O \texttt{SpeedGauge} converte velocidade em passos e movimenta o ponteiro físico. \\
\texttt{UC02} & Exibir rotação do motor no painel físico & O \texttt{RpmGauge} controla o instrumento de \textit{RPM} com suavização e resposta acelerada em variações bruscas. \\
\texttt{UC03} & Exibir marcha no display \texttt{TFT} & Os layouts do \texttt{DisplayService} apresentam a marcha de forma destacada no centro da interface. \\
\texttt{UC04} & Consultar status operacional do sistema & O estado \texttt{UiStatus} concentra conectividade, qualidade da telemetria, volta, posição e contadores de falha. \\
\bottomrule
\end{tabularx}
\end{table}

Casos de uso relacionados à compatibilidade multiplataforma e a efeitos
adicionais de vibração continuam pertencendo ao escopo analítico do TCC, mas
não foram integralmente convertidos em funcionalidade entregue no firmware
versionado.

\subsection{Diagrama de Classes}

A Figura~\ref{fig:diagrama_classes} apresenta a visão estrutural principal do
firmware. O diagrama deixa evidente a separação entre a classe orquestradora
\texttt{App}, os serviços de aplicação, os contratos de porta e os adaptadores
concretos de infraestrutura.

\FiguraRetratoGrande
{diagrama-de-classes.png}
{Diagrama de classes do firmware.}
{fig:diagrama_classes}

Do ponto de vista de responsabilidade, a classe \texttt{App} coordena o ciclo
de vida do sistema, o \texttt{TelemetryService} transforma entrada de rede em
estado de telemetria confiável, o \texttt{InstrumentCluster} despacha sinais
para instrumentos especializados e o \texttt{DisplayService} consolida a
renderização dos \textit{layouts} do painel. A modelagem também evidencia o uso de
interfaces como \texttt{IPacketReceiver}, \texttt{ITelemetryDecoder},
\texttt{IGaugeMotor}, \texttt{IStatusDisplay} e \texttt{IWifiConfigPortal},
decisão central para a testabilidade do projeto.

\subsection{Arquitetura do Sistema}

A arquitetura adotada combina organização em camadas com princípios de
\textit{Dependency Inversion} e \textit{Ports and Adapters}. O ponto de
composição do sistema localiza-se em \texttt{src/composition/main.cpp}, onde
adaptadores concretos, instrumentos, \textit{layouts} e serviços são instanciados e
injetados na aplicação.

\FiguraRetratoGrande
{arquitetura-geral.png}
{Arquitetura geral do sistema embarcado.}
{fig:arquitetura_geral}

Em termos de organização, a arquitetura pode ser resumida da seguinte forma:
\begin{itemize}
  \item \texttt{composition/}: monta a aplicação e concentra a criação das
  dependências concretas.
  \item \texttt{application/}: coordena fluxos de telemetria, instrumentos e
  interface sem depender diretamente de APIs de hardware.
  \item \texttt{domain/}: armazena modelos neutros, como
  \texttt{TelemetryFrame} e \texttt{UiStatus}.
  \item \texttt{ports/}: define contratos para tempo, rede, telemetria,
  display, botões e motores.
  \item \texttt{adapters/}: implementa \texttt{Wi-Fi}, \texttt{UDP},
  \texttt{display}, \texttt{GPIO}, relógio e acionamento dos instrumentos.
\end{itemize}

O fluxo principal de execução ocorre em quatro etapas: conexão de rede,
recepção de pacotes, decodificação/normalização da telemetria e despacho para
instrumentos e interface. No estado atual, a telemetria válida aciona
velocímetro, conta-giros, marcador de combustível e indicador de temperatura,
enquanto a interface exibe informações rápidas em quatro \textit{layouts} distintos.

\subsection{Diagramas de Sequência}

\FiguraRetrato
{diagrama-sequencia-inicializacao.png}
{Diagrama de sequência da inicialização do firmware.}
{fig:sequencia_inicializacao}

\FiguraRetrato
{diagrama-sequencia-operacao-principal.png}
{Diagrama de sequência da operação principal do firmware.}
{fig:sequencia_operacao_principal}

\FiguraRetrato
{diagrama-sequencia-telemetria.png}
{Diagrama de sequência de telemetria e instrumentos.}
{fig:sequencia_telemetria}

\subsection{Diagrama Entidade-Relacionamento}

Não se aplica ao escopo deste projeto. O firmware não utiliza banco de dados
relacional nem persistência tabular de domínio. O único armazenamento local
persistente implementado refere-se a credenciais de rede, mantidas pela API
\texttt{Preferences} do \texttt{ESP32}. Os dados de telemetria são mantidos em
memória volátil e atualizados em tempo real pelo \texttt{TelemetryService}.

\subsection{Interfaces}

A interface do sistema é inteiramente embarcada e se divide em dois contextos:
mensagens de estado operacional e layouts de telemetria do display
\texttt{TFT}. Como o repositório não versiona capturas de tela do painel físico,
esta seção descreve os \textit{layouts} concretamente implementados no código.

\begin{table}[H]
\centering
\caption{Interfaces visuais implementadas no display}
\label{tab:interfaces}
\footnotesize
\setlength{\tabcolsep}{4pt}
\renewcommand{\arraystretch}{1.25}
\begin{tabularx}{\textwidth}{
  >{\raggedright\arraybackslash}p{2.5cm}
  >{\justifying\arraybackslash}X
}
\toprule
\textbf{Interface} & \textbf{Conteúdo principal} \\
\midrule
\texttt{Layout0Basic} & Exibe posição de corrida, volta atual e marcha em destaque, privilegiando leitura rápida de estado competitivo. \\
\texttt{Layout1Minimal} & Exibe marcha central, velocidade, \textit{RPM} e percentual de combustível, servindo como painel sintético geral. \\
\texttt{Layout2Performance} & Exibe velocidade ampliada, marcha, combustível com codificação por cor e barra de \textit{RPM}. \\
\texttt{Layout3Tires} & Exibe temperaturas individuais \texttt{FL}, \texttt{FR}, \texttt{RL} e \texttt{RR} com codificação cromática por faixa térmica. \\
\texttt{Telas de estado} & Apresentam mensagens de conexão, orientação para acesso ao portal \texttt{SimHub} e aguardam configuração ou estabelecimento de rede. \\
\bottomrule
\end{tabularx}
\end{table}

A navegação entre \textit{layouts} é realizada por um botão físico ligado ao pino
\texttt{16}. O display é configurado em \texttt{240x320} pixels com driver
\texttt{ST7789}, conforme definido no arquivo \texttt{include/User\_Setup.h}.

\section{Software}
\label{sec:software}

\subsection{Implantação}

A implantação do projeto exige um ambiente de desenvolvimento compatível com
\texttt{PlatformIO}, uma placa \texttt{ESP32 Dev Module} e os componentes
físicos do protótipo, em especial o display \texttt{TFT}, os drivers
\texttt{TMC2208} e os instrumentos analógicos.

O processo de execução do firmware pode ser resumido nas etapas a seguir:
\begin{enumerate}
  \item Instalar ou disponibilizar o \texttt{PlatformIO CLI} e garantir que o
  projeto esteja aberto na raiz do repositório.
  \item Revisar o arquivo \texttt{include/User\_Setup.h} para confirmar a
  configuração do display \texttt{ST7789} e revisar
  \texttt{src/config/BoardConfig.h} para os pinos e limites dos instrumentos.
  \item Compilar o firmware principal com o comando
  \texttt{pio run -e esp32dev}.
  \item Transferir o binário para a placa por meio do fluxo padrão do
  \texttt{PlatformIO}.
  \item Ao energizar o sistema, observar se o dispositivo conecta-se a uma rede
  salva; caso isso não ocorra, conectar-se ao ponto de acesso \texttt{SimHub} e
  abrir \url{http://simhub} para informar SSID e senha da rede local.
  \item Configurar o simulador para enviar telemetria no formato atualmente
  suportado (\textit{Forza Motorsport 7 -- Data Out}) para a mesma rede local,
  utilizando a porta \texttt{5300}.
  \item Com a rede estabelecida, permitir que o firmware receba os pacotes e
  atualize tanto o display quanto os instrumentos físicos do painel.
\end{enumerate}

Para validação automatizada, os comandos principais do projeto são:
\begin{itemize}
  \item \texttt{pio test -e test\_native}
  \item \texttt{pio test -e test\_esp32\_smoke}
  \item \texttt{pio test -e test\_esp32\_udp\_integration}
  \item \texttt{./scripts/run\_native\_coverage.sh}
\end{itemize}

\subsection{Exemplo de Código}

O Código~\ref{cod:composition_root} apresenta um recorte do ponto de composição
do firmware, no qual as dependências concretas são instanciadas e conectadas à
classe principal da aplicação. Esse trecho é representativo porque evidencia,
em poucas linhas, a organização modular adotada no projeto.

\begin{lstlisting}[caption={Trecho do ponto de composição do firmware}, label={cod:composition_root}]
static ArduinoClock systemClock;
static UdpReceiver udp;
static Forza7Decoder decoder;
static TelemetryService telemetry(udp, decoder, systemClock);

static DisplayService ui(displayLayouts,
                         sizeof(displayLayouts) / sizeof(displayLayouts[0]));

static InstrumentCluster instruments(
  clusterInstruments,
  sizeof(clusterInstruments) / sizeof(clusterInstruments[0])
);

static App app(
  wifiPortal,
  layoutButton,
  systemClock,
  telemetry,
  ui,
  instruments
);

void setup() {
  app.begin(makeAppConfig());
}

void loop() {
  app.tick();
}
\end{lstlisting}

\subsection{Testes}

A estratégia de testes do projeto foi organizada em cinco camadas:
testes unitários, integração \textit{host-side}, teste de componente, \textit{smoke
test} embarcado e integração \texttt{UDP} real no \texttt{ESP32}. Essa divisão
foi adotada para permitir regressão rápida no host sem abrir mão de evidências
mínimas em hardware real.

\begin{table}[H]
\centering
\caption{Quadro consolidado de testes do projeto}
\label{tab:testes}
\footnotesize
\setlength{\tabcolsep}{4pt}
\renewcommand{\arraystretch}{1.25}
\begin{tabularx}{\textwidth}{
  >{\raggedright\arraybackslash}p{3.2cm}
  >{\centering\arraybackslash}p{1.3cm}
  >{\centering\arraybackslash}p{1.3cm}
  >{\justifying\arraybackslash}X
}
\toprule
\textbf{Tipo de teste} & \textbf{Suítes} & \textbf{Casos} & \textbf{Situação} \\
\midrule
Unitário & 12 & 51 & Implementado, executado e aprovado em \texttt{test\_native}. \\
Integração \textit{host-side} & 3 & 9 & Implementado, executado e aprovado em \texttt{test\_native}. \\
Componente & 1 & 8 & Implementado, executado e aprovado em \texttt{test\_native}. \\
Smoke embarcado & 1 & 2 & Implementado, executado e aprovado em \texttt{test\_esp32\_smoke}. \\
Integração embarcada real & 1 & 3 & Implementado; \textit{build} e \textit{upload} validados; execução final limitada pela bancada elétrica. \\
\bottomrule
\end{tabularx}
\end{table}

Os documentos de validação do repositório registram \texttt{18} suítes e
\texttt{73} casos definidos, dos quais \texttt{70} já tiveram execução concluída
com evidência coletada. A diferença entre o total documentado e a execução
\textit{host-side} decorre do fato de que \texttt{test\_native} cobre apenas as suítes
\textit{host-safe}, enquanto os demais casos pertencem aos ambientes embarcados.

Em verificação local realizada em 4 de maio de 2026, a suíte
\texttt{test\_native} foi executada novamente no estado atual do repositório e
produziu \texttt{68/68} casos aprovados em aproximadamente \texttt{13}
segundos. Esse resultado reafirma a integridade da automação \textit{host-side} na
versão final utilizada para a escrita deste relatório.

Quanto à cobertura real de código, a última medição consolidada registrada no
repositório aponta os percentuais apresentados na
Tabela~\ref{tab:cobertura}. Esses dados foram produzidos pelo script
\texttt{./scripts/run\_native\_coverage.sh}.

\begin{table}[H]
\centering
\caption{Resumo da cobertura \textit{host-side} registrada no projeto}
\label{tab:cobertura}
\small
\setlength{\tabcolsep}{6pt}
\renewcommand{\arraystretch}{1.2}
\begin{tabular}{lccc}
\toprule
\textbf{Grupo} & \textbf{Linhas} & \textbf{Funções} & \textbf{Branches} \\
\midrule
Unitário & 66,1\% & 79,5\% & 41,5\% \\
Integração & 13,4\% & 21,1\% & 6,6\% \\
Componente & 12,7\% & 15,3\% & 4,9\% \\
Consolidado \textit{host-side} & 76,9\% & 84,8\% & 46,2\% \\
\bottomrule
\end{tabular}
\end{table}

Do ponto de vista qualitativo, os testes cobrem os pontos tecnicamente mais
sensíveis do MVP: recepção \texttt{UDP}, \textit{offsets} do protocolo, conversões de
unidade, tratamento de \textit{timeout}, contagem de erros, filtros de entrada,
suavização de instrumentos, mudança de \textit{layout}, transições entre portal
\texttt{Wi-Fi} e modo operacional e renderização básica da interface. A
evidência embarcada já confirma o \textit{smoke test} no \texttt{ESP32}; por
outro lado, a integração \texttt{UDP} real permanece parcialmente limitada por
restrições físicas da bancada, especialmente condições de \textit{brownout}
quando o rádio \texttt{Wi-Fi} é ativado sob alimentação exclusivamente via
\texttt{USB}.

\section{Considerações Finais}

O repositório final do projeto demonstra que o TCC avançou de uma proposta de
painel háptico modular para um firmware embarcado efetivamente executável e
validável. A contribuição central da implementação não está apenas na recepção
de telemetria, mas na transformação dessa entrada em comportamento coerente de
um painel físico e visual, com separação arquitetural suficiente para manter o
código legível, testável e evolutivo.

Como resultado prático, o firmware entregue estabelece uma base funcional para
um painel de \textit{sim racing} com \texttt{ESP32}, portal de configuração
\texttt{Wi-Fi}, \textit{pipeline} de telemetria \texttt{UDP}, instrumentos físicos
calibrados e \textit{display} com \textit{layouts} especializados. A existência de testes
automatizados em várias camadas aumenta a confiança técnica no MVP e reduz o
risco de regressão em futuras iterações.

As limitações do trabalho também ficaram claras e devem ser registradas com
objetividade. A implementação atual concentra-se no protocolo
\textit{Forza Motorsport 7 (Data Out)}, não fecha ainda a compatibilidade
formal com múltiplas plataformas de jogo, não implementa os efeitos extras de
vibração previstos no escopo inicial e depende de melhor instrumentação de
bancada para concluir a validação \texttt{UDP} real em cenário eletricamente
estável. Assim, a evolução futura mais natural do projeto é consolidar a
validação experimental em uso real, medir latência fim a fim, ampliar o suporte
a novos simuladores e incorporar os efeitos hápticos adicionais previstos na
análise original.

% ------------------------------
\newpage
\renewcommand{\refname}{}
\section*{Referências}
\begin{thebibliography}{1}

\bibitem{wazlawick2017metodologia}
WAZLAWICK, Raul.
\textbf{Metodologia de pesquisa para ciência da computação}.
2. ed.
Rio de Janeiro: Elsevier Brasil, 2017.

\end{thebibliography}

% ------------------------------
\clearpage
\appendix
\section*{Apêndice A -- Manual de Configuração Rápida}
\addcontentsline{toc}{section}{Apêndice A -- Manual de Configuração Rápida}

\FiguraManual
{manual-de-uso.png}
{Manual de configuração rápida do painel de telemetria.}
{fig:manual_configuracao_rapida}

% ------------------------------
\clearpage
\section*{Anexos}
\addcontentsline{toc}{section}{Anexos}

\subsection*{Declaração de Entrega do Código-Fonte}
Declaro que o código-fonte correspondente a este trabalho foi entregue à
coordenação do curso e encontra-se organizado no repositório técnico utilizado
na elaboração deste relatório.

Assinatura: \rule{10cm}{0.4pt}

Data: \rule{1cm}{0.4pt}/\rule{1cm}{0.4pt}/\rule{2cm}{0.4pt}

\subsection*{Declaração de Submissão ao NIT}
Declaro que as providências institucionais relativas à submissão do software ao
NIT -- IFPI foram adotadas conforme as exigências aplicáveis ao presente
trabalho.

Assinatura: \rule{10cm}{0.4pt}

Data: \rule{1cm}{0.4pt}/\rule{1cm}{0.4pt}/\rule{2cm}{0.4pt}

\end{document}